\section{系统组成及功能说明}
\label{sec:methods}

\subsection{整体介绍}
\subsubsection{系统总体架构}
系统由图像采集、运动可视化、区域位移分析、视频显示和人机交互五个部分组成。Zynq处理系统（PS）负责设备配置与任务调度，可编程逻辑（PL）负责图像计算、视频接收及显示输出。图\ref{fig:system}给出了各模块的数据连接和控制关系。

\reportfig{fig2-1-system.pdf}{1.0}{系统总体架构与数据通路。实线表示图像或数值数据，虚线表示主要控制与状态连接；两条算法支路均读取原始采集帧。}{fig:system}

OV9281图像经CSI接收和像素整理后，由采集DMA写入DDR原始帧缓存。每幅完整原始帧供两个支路读取：全帧支路执行多尺度相位运动放大，生成增强图像；区域支路提取固定ROI，估计二维位移并计算时间频谱。显示模块根据所选视图读取原始帧或增强帧，叠加ROI框、频谱及运行状态，经HDMI输出。区域测量始终使用原始采集帧，使位移分析与运动增强显示分别对应原始运动和可视化结果。

PS通过控制寄存器设置帧地址和算法参数、发出启动请求，并根据采集、计算及显示的完成事件调度下一项任务。用户通过串口屏或实体按键操作系统，独立日志串口输出调试信息。计算模块产生图像和数值结果，PS程序汇总任务状态并向界面反馈执行结果。

\subsubsection{数据通路与软硬件分工}
以$1280\times800$、60帧每秒计算，视频输入规模约为每秒6144万个像素，处理过程还包含空间滤波、相位运算和图像重建。本设计将这些重复密集的运算组织为PL流水线，按数据访问方式配置行缓存、FIFO和运算单元；初始化、命令解析及异步事件处理由PS执行。Zynq在同一器件内集成处理器与可编程逻辑，为这种分工提供了硬件基础\cite{ug585}。表\ref{tab:partition}列出了各项任务的执行位置及数据接口。

\begin{table}[htbp]
\centering\small
\caption{主要任务的软硬件划分}\label{tab:partition}
\begin{tabularx}{\linewidth}{@{}p{2.6cm}p{1.3cm}X X@{}}
\toprule
任务 & 位置 & 输入与输出 & 划分依据 \\
\midrule
相机配置及采集控制 & PS & 配置请求至传感器状态反馈 & 以软件组织寄存器配置、启停和异常处理 \\
CSI接收及采集DMA & PL & RAW10数据至DDR完整帧 & 对连续输入进行像素整理和帧写入 \\
相位运动放大 & PL & 原始帧至处理帧 & 以流水线组织多尺度和相位运算 \\
ROI位移及时间频谱 & PL & 原始ROI至位移、谱幅 & 独立配置二维变换与窗口处理资源 \\
视频时序及叠加 & PL & 显示帧、频谱至HDMI & 在像素通路内生成显示内容 \\
命令、调度及状态确认 & PS & 用户请求至实际状态反馈 & 统一管理缓存所有权与生效条件 \\
\bottomrule
\end{tabularx}
\end{table}

系统采用分层存储组织：DDR保存完整图像及部分历史状态，AXI承担批量读写；FIFO、行缓存和BRAM工作区服务于局部算子。整体吞吐由算子处理速率、任务依赖、缓冲深度和DDR访问共同决定。FPGA在本设计中承担可定制计算通路及视频接口集成，视觉采集方式则支持对被测物进行非接触观察。

\subsection{硬件系统介绍}
\subsubsection{硬件整体介绍}
系统采用OS-ZYNQ7100开发板，主芯片为AMD Zynq-7100，BOM型号为\texttt{XC7Z100-2FFG900I}，Vivado目标器件为\texttt{xc7z100ffg900-2}，使用板载DDR和QSPI存储。外设包括OV9281摄像头、HDMI显示器、串口屏及供电装置，扬声器作为实验中的振动对象。外围硬件复用现有平台，核心设计集中于HLS计算IP、RTL视频逻辑和PS控制程序。

\subsubsection{机械设计介绍}
实验采用支架固定摄像头并对准扬声器，开发板和串口屏放置在便于操作的位置，具体布置见第三部分整机照片。稳定的相机视角有助于减少支架晃动、桌面运动及遮挡对图像的干扰。观察距离、视角、照明和固定方式是运动对照测试的控制条件，应在采集过程中保持一致。

\subsubsection{电路各模块介绍}
摄像头接口包括配置通路和图像数据通路：PS执行传感器寄存器配置，PL接收MIPI/CSI数据。DDR保存原始帧、处理帧及部分历史状态，显示DMA与视频时序逻辑完成HDMI输出。串口屏通过UART1和EMIO通信，UART0输出日志，板上按键经GPIO参与视图控制，QSPI存储启动镜像。接口引脚由工程约束文件指定，UART接线遵循板卡I/O电平及共地要求。芯片BOM、供电说明与启动拨码照片归档于\path{board/hardware/}，引脚和串口电平约定见\path{board/wiring/README.md}；串口屏HMI/TFT工程与下载说明由\path{board/hmi/README.md}索引。

\subsection{软件系统介绍}
\subsubsection{软件整体介绍}
控制程序运行于Arm Cortex-A9裸机环境，负责相机驱动、HMI协议、帧调度和状态反馈。HLS算法经综合形成在PL中执行的专用计算模块，RTL实现采集、显示和信息叠加逻辑。控制程序通过寄存器接口调用计算模块，数据流、流水线及存储映射由HLS综合和Vivado实现流程确定\cite{ug1399,ug901,ug904}。

上电后，程序初始化平台、计时器、通信接口和诊断状态，进入待机轮询。收到合法启动请求后，程序结束旧显示缓存的读取、配置相机并启动采集，随后装载算法参数并准备位移窗口。运行期间，主循环响应采集、计算及显示完成事件，处理用户命令和错误状态。停止时，程序先等待采集DMA在帧边界停止，再等待在途计算及频谱任务完成，最后关闭相机输出并反馈停止状态。图\ref{fig:main-flow}给出了各阶段的转换条件。

\reportfig{fig2-3-main-flow.pdf}{1.0}{主程序运行流程。启动就绪和停止排空分别触发运行、停止反馈；停止后显示可保持末帧，故障进入错误清理分支。}{fig:main-flow}

\subsubsection{软件各模块介绍}
\paragraph{（1）图像采集与帧缓存管理}
采集链路接收$1280\times800$的RAW10图像，像素以16位容器存入帧缓存，有效值范围为0至1023。程序检查完整帧及校验状态后提交算法处理，通过原始帧三缓冲协调采集和读取。调度器分别记录采集写入、运动放大、ROI计算及显示读取的完成事件。两条计算支路均释放原始帧后，该缓存才重新交给采集DMA，保证ROI任务读取期间的数据保持完整。

\paragraph{（2）多尺度相位运动放大}
相位运动放大利用局部相位随运动变化的关系增强细小运动的可见性\cite{wadhwa2013}。本设计将原始帧分解为$1280\times800$、$640\times400$和$320\times200$三级空间带通信号，并保留$160\times100$的粗尺度残差。各尺度依次完成局部相位计算、时间带通滤波和相位调制，再通过逐级上采样与叠加重建图像，如图\ref{fig:phase}所示。

\reportfig{fig2-4-phase.pdf}{1.0}{多尺度相位运动放大流程。各尺度独立保存参考相位与历史滤波状态，下部展开单尺度的处理关系。}{fig:phase}

对某一尺度中的一个像素，记带通信号为$B_t$，两个方向的空间响应为$R_{x,t}$、$R_{y,t}$，记两方向空间响应的幅值为$r_t$，本实现的局部相位定义为
\begin{equation}
 r_t=\sqrt{R_{x,t}^{2}+R_{y,t}^{2}}+\varepsilon,
 \qquad \phi_t=\operatorname{atan2}(B_t,r_t),
 \label{eq:local-phase}
\end{equation}
其中$\varepsilon=10^{-6}$避免幅度项退化，$t$为帧序号。以首帧相位$\phi_0$为参考，得到$x_t=\phi_t-\phi_0$。二阶IIR采用转置直接Ⅱ型形式，其更新关系为
\begin{equation}
\begin{aligned}
 y_t&=b_0x_t+z_{1,t-1},\\
 z_{1,t}&=-a_1y_t+z_{2,t-1},\\
 z_{2,t}&=b_2x_t-a_2y_t.
\end{aligned}
\label{eq:iir}
\end{equation}
上述关系描述实时处理模式，带通系数$b_1=0$。首帧设置参考相位并将两个IIR状态置零，因此$x_0=y_0=0$，首帧带通信号直接通过。默认频带为3至12\,Hz；PS以历史基线采样率常数$f_s=60.379777$\,Hz计算IIR系数，运行时按配置装载。以$\alpha$调节带通相位后，子带重建为
\begin{equation}
 B'_t=B_t\cos(\alpha y_t)+r_t\sin(\alpha y_t).
 \label{eq:band-reconstruction}
\end{equation}
重建图像经过范围钳位与取整后写入处理帧缓存。$\alpha$用于调节带通相位，其数值与物理位移放大倍率采用不同定义；实际运动变化还受输入运动及滤波响应影响。

空间邻域运算通过行缓存复用相邻像素，算子之间通过流式FIFO连接，FIFO深度用于容纳不同尺度的启动和处理节奏差异。参考相位及两个IIR状态均采用单精度浮点，每像素共96位有效状态，四个像素的状态紧凑存为三个128位AXI字，内部计算流仍采用128位表示；打包和解包模块完成布局转换。

\paragraph{（3）区域位移估计与频谱分析}
ROI从原始采集帧中截取，左上角固定为$(512,272)$，尺寸为$256\times256$。ROI计算与全帧放大并行执行，以首个有效ROI建立参考。当前ROI经二维Hann窗和二维FFT处理后，与参考频域数据计算归一化互功率谱。本工程采用自定义硬件实现，相位相关的基本方法见文献\cite{opencvphase}。

\reportfig{fig2-5-roi-spectrum.pdf}{1.0}{ROI位移估计与时间频谱流程。二维FFT/IFFT用于空间位移估计，直接DFT用于连续位移样本的时间频谱。}{fig:roi}

记参考和当前ROI的频域数据为$F_0$、$F_t$，对应实现可写成
\begin{equation}
 C_t=\frac{F_0\overline{F_t}}{|F_0\overline{F_t}|+10^{-9}},
 \qquad c_t=\mathcal{F}^{-1}\{C_t\}.
 \label{eq:phase-correlation}
\end{equation}
算法在相关面$c_t$中搜索峰值，沿$x$、$y$方向分别采用三点抛物线插值，得到位移$\mathrm{d}x$、$\mathrm{d}y$，单位为像素（px）。按式\eqref{eq:phase-correlation}的互功率谱顺序，理想平移时输出为将当前ROI对齐到参考ROI所需的位移，其符号与当前图像相对参考图像的平移方向相反。低纹理、像素取值异常或无效帧对应参考失效状态，失效样本按无效测量处理。

连续有效位移进入256点滑动窗口。时间频谱模块对每个方向去均值、去线性趋势并施加Hann窗，以直接DFT计算129个单边频点。记预处理样本为$\widetilde d[n]$，则
\begin{equation}
 D[k]=\sum_{n=0}^{255}\widetilde d[n]\exp\left(-\mathrm{j}\frac{2\pi kn}{256}\right),
 \quad f_k=\frac{kf_s}{256},\quad 0\leq k\leq128.
 \label{eq:temporal-dft}
\end{equation}
其中$f_s$为位移采样率。按历史实验中约60.38\,Hz的采样率计算，窗口时长$256/f_s$约4.24\,s，频点间隔$f_s/256$约0.236\,Hz。模块以约5\,Hz发布频谱，每次输出使用最近256个连续有效位移样本，相邻输出共享重叠样本。帧序号不连续时清空旧窗口，再从当前有效样本开始积累。频点间隔确定频率网格，频率估计偏差由第三部分的双频实验给出。

\begin{table}[htbp]
\centering\small
\caption{关键算法与显示参数}\label{tab:algorithm-parameters}
\begin{tabularx}{\linewidth}{@{}p{3.0cm}p{3.4cm}X@{}}
\toprule
参数 & 配置 & 设计含义 \\
\midrule
相位增益$\alpha$ & 0至50，默认12 & 调节带通相位的调制量 \\
时间带通 & 默认3至12\,Hz & 根据频带及参考采样率计算IIR系数 \\
ROI & $(512,272)$，$256\times256$ & 从原始采集帧截取的固定区域 \\
频谱窗口 & 256个连续位移样本 & 输出129个单边频点 \\
频谱发布 & 约5\,Hz & 控制频谱更新节奏，相邻输出窗口重叠 \\
HDMI输出 & $1920\times1080$，30\,Hz & 显示接口的分辨率与刷新时序 \\
\bottomrule
\end{tabularx}
\end{table}

\paragraph{（4）并行任务调度与结果提交}
调度程序向运动放大和ROI模块提交同一原始帧地址，分别维护两路在途任务和完成标记，并在两路均释放输入后回收缓存。通过状态检查的增强帧提交给显示模块，ROI位移结果则按帧序进入频谱窗口。参数变化或停止请求到达时，程序先处理在途任务，再确认新配置或结束排空，保证配置变更前后的输出对应正确参数。相机采集、处理帧提交及显示刷新分别计数，用于分析各环节的运行速率。

\paragraph{（5）视频显示与信息叠加}
系统提供原画、放大无面板和放大带频谱面板三种视图。ROI框标出位移分析区域，频谱面板显示两个方向的谱幅及运行信息。PL将图像、叠加内容与显示时序组合后输出。频谱数值结果独立保存，屏幕柱高按显示量程映射；图像位移谱幅采用像素单位。停止采集与计算后，显示可保留末帧，运行状态由控制程序反馈。

\paragraph{（6）人机交互与参数生效机制}
串口中断负责字节收发，主循环执行命令解析和业务操作。控制程序分别维护请求接收、配置排队、参数应用及状态反馈，使界面提示与设备执行阶段相对应。主要命令及生效条件见表\ref{tab:hmi-commands}。

\begin{table}[htbp]
\centering\small
\caption{交互命令及生效条件}\label{tab:hmi-commands}
\begin{tabularx}{\linewidth}{@{}p{2.2cm}p{2.2cm}X@{}}
\toprule
命令 & 参数范围 & 实际生效与反馈条件 \\
\midrule
\texttt{VIEW} & 0、1、2 & 选择三种视图；停止时切换视图仍保持停止状态 \\
\texttt{ALPHA} & 0至50 & 参数读回匹配且新参数处理帧完成后确认 \\
\texttt{EXPOSURE} & 1至898 & 空闲状态接收并排队，下次启动写入相机及读回；运行时返回忙状态 \\
\texttt{START} & 无 & 等待采集、处理及显示满足就绪条件 \\
\texttt{STOP} & 无 & 停止采集并排空相关任务，取消待执行重启后确认 \\
\bottomrule
\end{tabularx}
\end{table}

\reportfig{fig2-6-parameter-sequence.pdf}{1.0}{参数应用与状态确认时序。左图为运行中修改相位增益，右图为空闲时设置曝光并在下次启动应用；箭头和纵向排列表示操作先后关系。}{fig:parameters}

如图\ref{fig:parameters}所示，程序等待当前算法任务结束后写入$\alpha$，在寄存器读回匹配且新参数帧处理完成后反馈增益生效；曝光值在空闲时排队，下一次启动时写入相机并读回。\texttt{RUNNING}反馈要求显示已启动且至少完成两次显示；放大模式还要求至少完成一帧算法处理。协议采用设备状态反馈格式，消息中不设置请求序号。

\paragraph{（7）启动加载与工程复现}
启动链路由BootROM、FSBL、位流和应用组成。位流、XSA、平台初始化文件及ELF成套管理，使硬件地址、平台配置和应用寄存器布局保持对应。仓库按源码、IP、仿真、构建、板端记录及数据组织目录，提供主机测试、ARM编译、RTL比对、硬件构建和启动包生成入口。工程复现依次进行环境配置、构建、产物哈希核对及对应板卡的部署测试。

报告对应的2025.2硬件构建以归档导出IP为输入，HLS源码用于算法与接口追溯。11月3日补充的源码重生成已完成五个模块的C仿真、生产配置综合和IP导出，有限RTL覆盖范围见第三部分；新候选IP不替换本文资源表所对应的封存IP。源码到IP的重生成、等价回归，以及启动包静态检查、Flash烧写、冷启动和业务运行，分别对应不同的复现步骤。各版本的测试对象和完成结果在第三部分列明，具体入口见第\ref{sec:reproduction}节。
